\chapter{Root cause as an inverse problem}
\label{ch:inverse}

\section{The forward map and its ill-posedness}
Unrolling \eqref{eq:propagation} from onset, the degradation field is a linear function of the
injection history,
\begin{equation}
d(t)=\sum_{k=0}^{t-t_0}M^{k}\,\mathbf 1_{v_0}\,s(t-k),\qquad M=\rho I+\beta P,\quad P=D_{\mathrm{in}}^{-1}A^{\top}.
\label{eq:forward}
\end{equation}
The root-cause question is: given noisy observations of $X$ (not $d$), recover $v_0$. This is an
inverse problem on the graph and it is ill-posed in the usual ways. Two origins with the same
downstream reach are indistinguishable from their effects alone (a link and the router it feeds
differ only in the link's own KPIs). The observation \eqref{eq:observation} mixes $d$ with diurnal
and coloured noise. And $M^k$ smooths: the impulse response spreads over $k$ hops, so late
observations carry less positional information than early ones.

What the structure of \eqref{eq:forward} tells us is which inductive biases a learned inverse must
contain. The adjoint of the forward operator is $M^\top=\rho I+\beta P^\top$, which propagates
information effect$\to$cause along the transposed relations. That is why $\Rel^{\top}$ is in
\eqref{eq:rgcn}: a network without transposes can only smooth downstream and cannot localise.

\section{Learned inverse}
We do not invert \eqref{eq:forward} analytically; real networks are not exactly linear, the
observation noise is not exactly Gaussian, and the ramp is unknown. We learn a conditional
distribution over origins from the relational embedding. With the logit $z_v(t)$ from
\eqref{eq:heads},
\begin{equation}
p(v_0=v\mid X_{[t-W,t)},G)=\frac{\exp z_v(t)}{\sum_{u\in\V}\exp z_u(t)},
\label{eq:rca}
\end{equation}
and the root-cause ranking at time $t$ is $\V$ sorted by $z_v(t)$ descending. The softmax is over
\emph{all nodes of all types}: the network is asked ``which single node'' across RAN, Core and
Transport, which is the cross-domain triage question.

\begin{remark}[Node scoring is link prediction with one anomaly node]
\label{rem:linkpred}
A common production formulation adds an \emph{anomaly} node $a$ to the property graph, connected
to every symptomatic entity, and trains a link predictor to score the edge $a\to v$ labelled
``originates from''. With a bilinear scorer $z_v=h_a^{\top}Mh_v$ and a single anomaly node per
incident, $h_a$ is a constant for the episode and $z_v$ is a linear read-out of $h_v$: exactly
\eqref{eq:heads}. The two formulations differ only when several concurrent incidents must be
attributed separately, in which case the anomaly-node version is the natural generalisation.
\end{remark}

\begin{remark}[Equivariance makes the ranking well defined]
Because $z$ is a per-node output of an equivariant map \eqref{eq:equivariance}, relabelling the
network permutes the ranking consistently. The claim ``node 17 is the root cause'' is really the
claim ``the node with this history and this neighbourhood is the root cause'', which is what an
operator means.
\end{remark}

\section{Loss}
The three heads are trained jointly on windows drawn from simulated episodes,
\begin{equation}
\mathcal L=\underbrace{\frac1{NK}\norm{\hat X(t)-X(t)}_F^2}_{\text{forecast}}
+\underbrace{\frac1N\norm{\hat d(t)-d(t)}_2^2}_{\text{degradation}}
+\lambda\,\underbrace{\ind{t\ge t_0+\delta}\big(-\log p(v_0\mid\cdot)\big)}_{\text{root cause}},
\label{eq:loss}
\end{equation}
with the root-cause term switched on $\delta$ steps after onset, when some signal has
accumulated, and switched off entirely on nominal episodes. The forecast term is included even
though detection uses $\hat d$: it is label-free, it regularises the embedding toward the KPI
dynamics, and it is the term that survives when a deployment has no incident labels.

\section{Counterfactual check}
Because the twin is explicit, a ranking can be verified rather than trusted: remove the injection
at the top-ranked node and re-simulate. If the field decays, the ranking was right. The closed loop
of Chapter~\ref{ch:closed} does exactly this, and the exposure numbers in
Chapter~\ref{ch:experiments} are computed from the re-simulation, not from the model's own
confidence.
